\section{The pair-profile estimate}\label{sec:two}
We prove Proposition~\ref{prop:pair} by showing that two profiles with $\mu[c_i]<\Lambda_0$ and total half-line mass at most $1/2$ must have excess above $0.9$ greater than $0.0872$.

For positive bounded profiles on a finite interval, use $t=\sqrt{\Lambda_0}\,y$, $\ell_i=\sqrt{\Lambda_0}\,c_i$ and the scaled Pr\"ufer angles
\[
 \theta_i'=1-\frac{\pi^2\sin^2\theta_i}{\ell_i^2},\qquad \theta_i(0)=\pi/2.
\]
Sturm theory~\cite{Zettl} says that $\mu[c_i]<\Lambda_0$ forces each angle to reach $\pi$ in finite time. General profiles follow by truncating to an interval containing compactly supported test functions, capping large widths, and adding a small positive width. The strict Rayleigh inequalities survive these changes; the extra mass and excess tend to zero. The certificate below has strict slack in both the initial state and final excess, so this approximation does not lose the conclusion.

Set $C=3\sqrt3\pi/2$, $U_i=1+\cos\theta_i$ and $X_i=\sin\theta_i$. Before the first angle reaches $\pi$, $0<\theta_i<\pi$, so $X_i\ge0$. The elementary inequality
\begin{equation}\label{eq:angle-budget}
 \ell_i\ge C X_i\theta_i'
\end{equation}
follows, on writing $\ell_i=\sqrt3\pi z$, from $2z+X_i^3/z^2\ge3X_i$; the difference is $(z-X_i)^2(2z+X_i)/z^2$.

Let $M_*$ be the allowed physical half-line mass, initially $1/2$, and let $M(t)$ be the mass already used. Introduce
\[
 s(t)=\frac{\Lambda_0}{C}\bigl(M_*-M(t)\bigr)-U_1(t)-U_2(t).
\]
By~\eqref{eq:angle-budget}, $s$ is nonincreasing. It is nonnegative along a feasible trajectory: applying~\eqref{eq:angle-budget} to each angle's remaining journey to $\pi$ bounds the remaining mass below by $C(U_1+U_2)/\Lambda_0$. At the start,
\[
 s(0)=\frac{\Lambda_0}{3\sqrt3\pi}-2.
\]
Thus the trajectory lies in $[0,2]^2\times[0,s_{\rm top}]$, where the certificate chooses $s_{\rm top}$ slightly larger than this initial value.

Let $w(U_1,U_2,s)$ be continuous and affine on each tetrahedron of a triangulation. Write its gradient on a tetrahedron as $(q_1,q_2,-k)$, and put $P_i=q_i+k$. Require $k,P_1,P_2\ge0$ and $w\le0$ on each face $U_i=0$. Put
\[
 \omega=\frac{\sqrt{\Lambda_0}\,0.9}{\sqrt3\pi},\qquad
 A=P_1^{1/3}X_1+P_2^{1/3}X_2.
\]
The remaining cell condition is
\begin{equation}\label{eq:HJB}
 3(P_1X_1+P_2X_2)\le\Psi_k(A),
\end{equation}
where
\[
\Psi_k(A)=
\begin{cases}
3k^{2/3}A,&A\le k^{1/3}\omega,\\
2k\omega+A^3/\omega^2,&k^{1/3}\omega\le A\le(k+1/2)^{1/3}\omega,\\
3(k+1/2)^{2/3}A-\omega,&A\ge(k+1/2)^{1/3}\omega.
\end{cases}
\]
To see why this is sufficient, write $\ell_i=\sqrt3\pi z_i$ and $Z=z_1+z_2$. Direct differentiation shows that the required inequality along a trajectory is
\[
 \sum_i P_i\left(3X_i-\frac{X_i^3}{z_i^2}\right)
 \le 2kZ+(Z-\omega)_+.
\]
At fixed $Z$, H\"older's inequality gives
\[
 \inf_{z_1+z_2=Z}\sum_i\frac{P_iX_i^3}{z_i^2}=\frac{A^3}{Z^2}.
\]
Minimizing $A^3/Z^2+2kZ+(Z-\omega)_+$ over $Z>0$ gives exactly the three branches of $\Psi_k$. Thus~\eqref{eq:HJB} proves, for every control,
\[
 \frac{\dd}{\dd t}\left(\frac{2C}{\Lambda_0}w\right)
 \ge-\frac{(\ell_1+\ell_2-\sqrt{\Lambda_0}\,0.9)_+}{\Lambda_0}.
\]
The chain rule also holds when the trajectory follows a cell face: continuity of the affine pieces makes their tangential derivatives agree.

Stop when the first angle reaches $\pi$. The terminal value of $w$ is nonpositive. Since $w_s=-k\le0$, its initial value is at least $w(1,1,s_{\rm top})$. Integration gives
\[
 \int_0^\infty(c_1+c_2-0.9)_+\dd y
 \ge\frac{3\sqrt3\pi}{\Lambda_0}w(1,1,s_{\rm top}).
\]
The rational checker verifies that this lower bound exceeds $0.0879293$. This is strictly larger than $0.0872$, proving Proposition~\ref{prop:pair}.
